% TOP_PROBLEM: {"problemId": "problem.kaplansky-s-zero-divisor-conjecture-for-torsion-free-groups", "canonicalTitle": "Kaplansky's Zero-Divisor Conjecture for Torsion-Free Groups", "releaseRank": 204, "primaryDomain": "algebra_representation_category", "primaryDomainLabel": "Algebra, representation & category theory", "displayStatus": "open", "releaseStatus": "open"}
% !TeX program = lualatex
% Definition and sources only; this file is not an LLM solution attempt.
\documentclass[11pt]{article}
\usepackage[margin=25mm]{geometry}
\usepackage{fontspec}
\setmainfont{Latin Modern Roman}
\usepackage{amsmath,amssymb,mathtools}
\usepackage{unicode-math}
\setmathfont{Latin Modern Math}
\usepackage{xurl,hyperref}
\hypersetup{colorlinks=true,urlcolor=blue}
\setcounter{secnumdepth}{0}
\setlength{\parindent}{0pt}
\setlength{\parskip}{5pt}
\setlength{\emergencystretch}{3em}
\begin{document}
\section*{\texorpdfstring{Kaplansky's Zero-Divisor Conjecture for Torsion-Free Groups}{Kaplansky's Zero-Divisor Conjecture for Torsion-Free Groups}}\label{204}
\subsection{Definitions and mathematical statement}
For a field $K$ and a group $G$, the group ring $K[G]$ consists of finite sums $\sum_{g\in G}a_gg$, with coefficientwise addition and multiplication determined by the group law and distributivity. A group is torsion-free if $g^m=e$ with $m\ge1$ implies $g=e$. The conjecture is
\[
 a,b\in K[G],\quad ab=0\quad\Longrightarrow\quad a=0\text{ or }b=0
\]
for every torsion-free $G$ and every field $K$. The ring may be noncommutative; ``domain'' here means the displayed absence of nonzero zero-divisors, not commutativity. No finite generation or restriction on the characteristic is imposed.
\par\smallskip\textit{Formulation and concept references:} \href{https://arxiv.org/abs/2501.07646}{[S1]}.

\subsection{Short English statement}
Can two nonzero finite linear combinations of elements of a torsion-free group ever multiply to zero over a field?
\subsection{Sources}
\begin{itemize}
\item[S1]M. Garg, I. Mineyev, arXiv:2501.07646 (2025).
\url{https://arxiv.org/abs/2501.07646}.
\textit{Formulation source.}
\end{itemize}
\end{document}
